%!TEX root = ../../thesis.tex

\section{Rewards}\label{sec:mdp-rewards}

Executing the \gls{sut} produces a list of block hit counts from which the reward is built.
We do not use the hit counts directly, but rather map them to an \gls{afl}-style hit count \emph{level}\index{hit count level}

\begin{equation*}
  \ell(n)=
  \begin{cases}
    0 & n=0, \\
    \min\left(1+\lfloor\log_2 n\rfloor,\ \ell_{\max}\right) & n>0,
  \end{cases}
\end{equation*}

and a \emph{coverage score}\index{coverage score} is derived from the hit count levels by summing over all buckets
$\sigma(c)=\sum_i\ell(c_i)$ of an execution $c\in\mathcal{C}$.
At 4096 buckets and $\ell_{\max}=8$, a single execution scores between 0 and 32768.
Compressing the raw counts into exponentially spaced buckets makes the coverage score less sensitive
to the exact hit count and represents a wider range of hit counts than linear
bucketing does.

\subsection{Input-level actions}

In the input-level actions setting, the fuzzer executes single inputs and keeps no execution history around.
Rewards are designed in a simple and straightforward manner that serves the purpose of the target.
For example, the sequence (\cref{sec:target-sequence}) example receives a fixed liveliness reward for
each step it keeps the environment alive, as well as a final coverage reward at the end that
includes additional bonus reward if the password is found.
The high-and-low (\cref{sec:target-high-and-low}) example is awarded the raw sum of block hit counts.

\subsection{Corpus-level actions}

In general, rewards for corpus-level actions can either be \emph{memoryless}
or \emph{marginal}.
Memoryless rewards\index{reward!memoryless} depend only on the coverage score of each execution and ignore
any coverage reached by previous seeds.
In contrast, marginal coverage\index{reward!marginal} pays only for the levels that have not been reached before.
\Textcite{bottinger_deep_2018} evaluate both and report that only the memoryless variant
beats a random baseline, while the variant built on coverage history does not.

The memoryless variant can be interpreted as steering the seed mutation towards high-coverage
regions.
The problem is that mutations produce the same high-coverage seed repeatedly.
Paying out the same high-coverage reward for seeds that already reach high coverage
discourages exploration into new regions.

We solve this problem by introducing a marginal reward.
By keeping track of the union of all coverage seen during an episode, we can
subtract this union from a new coverage map to get the marginal coverage.
This contains only newly found coverage and thus encourages exploration into yet underexplored regions.
A disadvantage of marginal coverage is its time dependency.
A seed that achieves new coverage early in the episode when the coverage union is still sparse
achieves a higher marginal reward than the same seed later on in the run, even though
the later seed might produce a new coverage high score.
We therefore include a depth term that rewards executions with a coverage score
larger than that of any previous execution.

An episode maintains two running coverage maxima over its executions: the
\emph{union coverage}\index{union coverage} and the \emph{depth record}\index{depth record}.
Each bucket $i$ of the union coverage $u_t$ contains the maximum of the $i$th bucket of all previous
executions, $u_{t,i}=\max_{k<t}c_{k,i}$, while the depth record $d_t=\max_{k<t}\sigma(c_k)$
represents the highest coverage score a single execution has produced.

The reward of an execution $c_t$ is then expressed as the combination of a marginal term
that encourages exploration breadth by rewarding new levels and of a depth term
that encourages exploration depth by rewarding \emph{deep} executions
that improve on the previous executions with the highest coverage score.

\begin{equation*}
  r_t=\lambda\left[\sum_i\bigl(\ell(c_{t,i})-\ell(u_{t,i})\bigr)_+
  \;+\;\omega\bigl(\sigma(c_t)-d_t\bigr)_+\right],
\end{equation*}

where $\omega$ is a parameter that adjusts the weight of exploration depth against
breadth, and $\lambda$ rescales the reward to better fit the value support.
For the default choice of $\omega=1$, we end up with a reward signal that
values an execution that reaches exactly one level deeper in a single bucket
than the union of all previous executions and an execution that reaches
exactly one level deeper than the best previous execution equally.
For $\omega=0$, the depth term disappears and the reward collapses to a pure measure
of marginal novelty, i.e., of how much a new seed improves on the union of all previously attained coverage.
